\documentclass[conference]{IEEEtran}
\usepackage{graphicx}
\usepackage{amsmath}
\usepackage{multirow}
\usepackage{booktabs}
\usepackage{makecell}
\usepackage{tabularx}
\usepackage[numbers]{natbib}
\usepackage{multicol}
\usepackage[bookmarks=true]{hyperref}
\usepackage{cuted}
\usepackage{caption}
\usepackage{algorithm}
\usepackage{algorithmic}
\usepackage{xcolor}
\bibliographystyle{unsrt}

% \setlength{\textfloatsep}{2pt plus 2pt minus 2pt}
% \setlength{\intextsep}{2pt plus 2pt minus 2pt}
% \setlength{\abovecaptionskip}{2pt}
% \setlength{\belowcaptionskip}{2pt}

\setlength{\textfloatsep}{1pt plus 1pt minus 1pt}
\setlength{\intextsep}{1pt plus 1pt minus 1pt}
\setlength{\abovecaptionskip}{1pt}
\setlength{\belowcaptionskip}{1pt}

\pdfinfo{
   /Author (Author Names Omitted)
   /Title  (MIF-Humanoid Rebuttal)
   /Subject (Robotics)
}

\begin{document}


\twocolumn[
\begin{center}
    \textbf{\Large Rebuttal for Paper-ID 256}
\end{center}
\vspace{-5pt}
]

The authors thank the Area Chair and the reviewers for constructive feedback. Below we address the key concerns.

% \textbf{1) Contribution (novelty/positioning).}
% Our claim is not that we achieve substantial novel improvements in each of the semantic 3DGS, scene graphs, feature compression, or object-level reconstruction. The contribution is their \emph{humanoid-oriented integration} for navigation:
% (1) a confidence-gated appearance field that reduces gait-induced artifacts during online semantic 3DGS mapping,
% (2) a discrepancy-triggered spatial field that updates only locally inconsistent regions under scene changes,
% (3) a geometry field that performs task-driven active view selection for safer interaction, and
% (4) we validate the integrated system on a real humanoid platform in dynamic environments. 
\textbf{1) Contribution (novelty/positioning).}
Our contribution is the \emph{humanoid-oriented integration} of these components for the new manipulation-safety awared navigation task: (1) confidence-gated appearance field reducing gait artifacts in online semantic 3DGS, (2) discrepancy-triggered spatial field updating only inconsistent regions, (3) geometry field with manipulation-driven active view selection and interaction pose selection, and (4) the system is deployable according to memory usage and latency, and we do validation on a real humanoid in dynamic environments.




\textbf{2) Additional evidence on the main disputed points.}

% \textbf{Other 3DGS baselines and confidence formulation.}
% We added comparisons with Deblur-GS and PUP-3DGS when the robot is operated in slow and fast-walk conditions, as shown in Table~\ref{tab:conf}.
% $C_i$ is computed from each Gaussian's normalized optimization instability (weighted gradient magnitude of primitive attributes, including 3D position, scale/rotation, color, and opacity): gait-corrupted primitives remain persistently high-gradient across training, while static textured boundaries converge, so low-$C_i$ primitives are downweighted in rendering and object reliability estimation. The primary noise source targeted by $C_i$ is locomotion instability from humanoid motion, as well as local pose estimation errors in SLAM.

\textbf{Other 3DGS baselines and confidence formulation.}
We added comparisons with Deblur-GS and PUP-3DGS under slow-walk and fast-walk conditions (Table~\ref{tab:conf}). Different from these methods, $C_i$ is computed from each Gaussian's gradient magnitude (weighted across position, scale/rotation, color) and opacity, both normalized by global means.
 Gait-corrupted primitives exhibit persistently high gradients across training, while static high-texture boundaries converge naturally. Leveraging this, $C_i$ modulates both mapping and tracking: low-confidence primitives receive higher learning rates for refinement in mapping, while being masked in photometric pose optimization to prevent noise propagation. The primary noise sources targeted by $C_i$ are locomotion instability from humanoid motion and local pose errors in SLAM.


% Rendered views from the appearance field provide more robust object detection (YOLO + SAM) and spatial relation extraction (GPT-4o) by reducing gait-induced artifacts for better scene graph construction.(Table~\ref{tab:render}).


\begin{table}[h]
\centering
\caption{ Rendering results under slow (S-) and fast (F-) walk.}
\label{tab:conf}
% 稍微收紧列间距，防止等比缩放时字体变得过小
\setlength{\tabcolsep}{2.5pt} 
% 自动将表格宽度缩放至当前单栏的宽度 (\columnwidth)
\resizebox{\columnwidth}{!}{%
\begin{tabular}{lccccc}
\toprule
\textbf{Method} & \textbf{S-PSNR$\uparrow$} & \textbf{S-SSIM$\uparrow$} & \textbf{F-PSNR$\uparrow$} & \textbf{F-SSIM$\uparrow$}\\
\midrule
w/o Confidence & 28.4 & 0.88 & 21.6 & 0.72\\
Deblur-GS [I3D'24] & 25.14 & 0.80 & 24.67 & 0.70\\
PUP-3DGS [CVPR'25] & 20.44 & 0.67 & 17.89 & 0.53\\
\textbf{MIF $C_i$ (Ours)} & \textbf{31.2} & \textbf{0.93} & \textbf{29.8} & \textbf{0.89}\\
\bottomrule
\end{tabular}%
}
\end{table}

The confidence module is also robust to hyperparameter choice: over $\beta\!\in\![2,8]$ and $\gamma\!\in\![0.5,4]$, PSNR varies by less than 0.5\,dB and FPR by less than 3\%.





\textbf{Dynamic-scene baselines.}
Our original comparison to HOV-SG was intended to expose the failure mode of static-memory methods under scene changes, not to claim superiority over all dynamic scene graph approaches. We agree that static baselines alone are insufficient and therefore added comparisons with Khronos and a periodic-rescan ConceptGraphs variant in the same non-revisiting protocol. As shown in Table~\ref{tab:dsg}, we test whether each method can update the scene graph (against changes within the environment) based on the data collected during navigation after the environment change.

\begin{table}[h]
\centering
\caption{Change detection under non-revisiting navigation. Reloc/Remov/Addit: object relocated ($>$1.5m), removed, or added out-of-view. Robot navigates directly without revisiting.}
\label{tab:dsg}
\footnotesize
\setlength{\tabcolsep}{2pt}
\begin{tabularx}{\columnwidth}{@{} l
>{\centering\arraybackslash}X
>{\centering\arraybackslash}X
>{\centering\arraybackslash}X @{}}
\toprule
\textbf{Method} & \textbf{Reloc SR$\uparrow$} & \textbf{Remov SR$\uparrow$} & \textbf{Addit SR$\uparrow$}\\
\midrule
ConceptGraphs + Rescan & 38\% & 35\% & 12\% \\
Khronos [RSS'24] & $\sim$18\% & 61\% & $\sim$12\% \\
\textbf{MIF (Ours)} & \textbf{94\%} & \textbf{98\%} & \textbf{86\%} \\
\bottomrule
\end{tabularx}
{\tiny\raggedright $^*$Two-session protocol: initial mapping, then scene changes, followed by change detection via scene graph comparison. The robot during navigation in two session does not ensure exactly the same trajectory.\par}
\vspace{-3mm}
\end{table}

\textbf{$\mathcal{D}$-Score \& threshold selection.} Under gait noise, $\mathcal{D}$ remains concentrated in a low-value, low-variance regime (mean $\sim$0.2, std $\sim$0.05), while real changes form a separable high-value mode (mean $\sim$0.6, std $\sim$0.1), enabling reliable discrimination. We also report the ROC-based selection of change trigger threshold $\tau$ in Table~\ref{tab:tau}; $\tau=0.45$ maximizes F1 and was selected from this quantitative evaluation rather than manual tuning.

\begin{table}[h]
\centering
\caption{ROC-based threshold selection in real scenarios.}
\label{tab:tau}
\footnotesize
\begin{tabularx}{\columnwidth}{@{}
>{\centering\arraybackslash}X
>{\centering\arraybackslash}X
>{\centering\arraybackslash}X
>{\centering\arraybackslash}X @{}}
\toprule
\textbf{Threshold $\bm{\tau}$} & \textbf{True Positive Rate$\uparrow$} & \textbf{False Positive Rate$\downarrow$} & \textbf{F1$\uparrow$}\\
\midrule
0.25 & 99.2\% & 18.7\% & 0.817\\
0.35 & 97.8\% & 9.3\% & 0.893\\
\textbf{0.45 (Ours)} & \textbf{95.1\%} & \textbf{4.2\%} & \textbf{0.934}\\
0.55 & 89.6\% & 1.8\% & 0.913\\
0.65 & 78.3\% & 0.7\% & 0.866\\
\bottomrule
\end{tabularx}
\end{table}

\textbf{3) Other clarifications requested by reviewers.}

\textbf{Viewpoint utility $Q(k)$.}
We agree the term was imprecise. $Q(k)$ is not information-theoretic but a \emph{task-local geometric observation utility} for manipulation safety. Candidate views are sampled on a hemisphere around the object, and $Q(k)$ ranks them by distance, foveal alignment, and high-confidence Gaussian density to maximize interaction surface observability.

\textbf{Semantic compression.}
Semantic compression is not claimed as a standalone novelty. It is an efficiency-enabling choice that preserves task-relevant semantics while making online deployment practical, as shown in Table~\ref{tab:semantic_embedding}.

\begin{table}[h]
\centering
\caption{Semantic localization under slow (S-) and fast (F-) walk conditions, memory usage, and latency.}
\label{tab:semantic_embedding}
\setlength{\tabcolsep}{3pt} % 收紧列间距
\resizebox{\columnwidth}{!}{% 自动缩放至单栏宽度
\begin{tabular}{lcccc}
\toprule
\textbf{Method} & \textbf{VRAM$\downarrow$\,(GB)} & \textbf{Latency$\downarrow$\,(s)} & \textbf{S-SR$\uparrow$} & \textbf{F-SR$\uparrow$}\\
\midrule
LangSplat [CVPR'24] & 18.45 & 2.3 & 95\% & 72\%\\
LatentBKI-64D [RA-L'25] & 23.5 & 7.1 & \textbf{98\%} & 74\%\\
FeatSplat-32D [ECCV'24] & 4.8 & 3.2 & 81\% & 59\%\\
\textbf{MIF 64D (Ours)} & 18.0 & 4.6 & 92\% & \textbf{88\%}\\
\textbf{MIF 32D (Ours)} & \textbf{$<$4.0} & \textbf{1.6} & 90\% & 86\%\\
\bottomrule
\end{tabular}}
\end{table}

\textbf{Rendered views for the spatial field.}
The appearance field renders stabilized views that reduce gait-induced artifacts, enabling more robust object detection (YOLO + SAM) and spatial relation extraction (GPT-4o) for scene graph construction (Table~\ref{tab:render}). The discrepancy score compares the node position, semantics, and relation sets of the local and global graph to trigger local updates when $\mathcal{D}>\tau$.

\begin{table}[h]
\centering
\caption{Semantic object localization on scene graphs constructed from rendered vs. raw frames.}

\label{tab:render}
\setlength{\tabcolsep}{4pt} % 稍微收紧列间距，确保第一列不折行
\footnotesize % 严格保持和之前表格完全一致的字号
% 彻底移除 resizebox，使用 tabularx 强制总宽度为 \columnwidth
% 第一列用 l (左对齐)，后三列用 X (居中且自动均分剩余宽度)
% @{} 用于去除左右两端多余空白，确保边缘完美贴合题注
\begin{tabularx}{\columnwidth}{@{} l 
    >{\centering\arraybackslash}X 
    >{\centering\arraybackslash}X 
    >{\centering\arraybackslash}X @{}}
\toprule
\textbf{Input Source} & \textbf{Recall$\uparrow$} & \textbf{Precision$\uparrow$} & \textbf{mAP$\uparrow$}\\
\midrule
Raw frames (Slow) & 0.87 & 0.91 & 0.84\\
Raw frames (Fast) & 0.61 & 0.79 & 0.58\\
\textbf{MIF Rendered (Slow)} & \textbf{0.94} & \textbf{0.95} & \textbf{0.93}\\
\textbf{MIF Rendered (Fast)} & \textbf{0.92} & \textbf{0.94} & \textbf{0.91}\\
Static camera (Ref.) & $\sim$1.00 & $\sim$1.00 & $\sim$1.00\\
\bottomrule
\end{tabularx}

\end{table}

\textbf{Others.}
Moreover, we will:
(1) expand related work on dynamic scene graphs, confidence in 3DGS, motion-deblurring 3DGS, and semantic feature distillation and
(2) clarify all notations and variables.
(3) release the codes for both the algorithm and the whole ROS system upon acceptance. 
\end{document}